% Options for packages loaded elsewhere
% Options for packages loaded elsewhere
\PassOptionsToPackage{unicode}{hyperref}
\PassOptionsToPackage{hyphens}{url}
\PassOptionsToPackage{dvipsnames,svgnames,x11names}{xcolor}
%
\documentclass[
  english,
  a4paper,
  12pt]{article}
\usepackage{xcolor}
\usepackage{amsmath,amssymb}
\setcounter{secnumdepth}{5}
\usepackage{iftex}
\ifPDFTeX
  \usepackage[T1]{fontenc}
  \usepackage[utf8]{inputenc}
  \usepackage{textcomp} % provide euro and other symbols
\else % if luatex or xetex
  \usepackage{unicode-math} % this also loads fontspec
  \defaultfontfeatures{Scale=MatchLowercase}
  \defaultfontfeatures[\rmfamily]{Ligatures=TeX,Scale=1}
\fi
\usepackage{lmodern}
\ifPDFTeX\else
  % xetex/luatex font selection
\fi
% Use upquote if available, for straight quotes in verbatim environments
\IfFileExists{upquote.sty}{\usepackage{upquote}}{}
\IfFileExists{microtype.sty}{% use microtype if available
  \usepackage[]{microtype}
  \UseMicrotypeSet[protrusion]{basicmath} % disable protrusion for tt fonts
}{}
\makeatletter
\@ifundefined{KOMAClassName}{% if non-KOMA class
  \IfFileExists{parskip.sty}{%
    \usepackage{parskip}
  }{% else
    \setlength{\parindent}{0pt}
    \setlength{\parskip}{6pt plus 2pt minus 1pt}}
}{% if KOMA class
  \KOMAoptions{parskip=half}}
\makeatother
% Make \paragraph and \subparagraph free-standing
\makeatletter
\ifx\paragraph\undefined\else
  \let\oldparagraph\paragraph
  \renewcommand{\paragraph}{
    \@ifstar
      \xxxParagraphStar
      \xxxParagraphNoStar
  }
  \newcommand{\xxxParagraphStar}[1]{\oldparagraph*{#1}\mbox{}}
  \newcommand{\xxxParagraphNoStar}[1]{\oldparagraph{#1}\mbox{}}
\fi
\ifx\subparagraph\undefined\else
  \let\oldsubparagraph\subparagraph
  \renewcommand{\subparagraph}{
    \@ifstar
      \xxxSubParagraphStar
      \xxxSubParagraphNoStar
  }
  \newcommand{\xxxSubParagraphStar}[1]{\oldsubparagraph*{#1}\mbox{}}
  \newcommand{\xxxSubParagraphNoStar}[1]{\oldsubparagraph{#1}\mbox{}}
\fi
\makeatother


\usepackage{longtable,booktabs,array}
\usepackage{calc} % for calculating minipage widths
% Correct order of tables after \paragraph or \subparagraph
\usepackage{etoolbox}
\makeatletter
\patchcmd\longtable{\par}{\if@noskipsec\mbox{}\fi\par}{}{}
\makeatother
% Allow footnotes in longtable head/foot
\IfFileExists{footnotehyper.sty}{\usepackage{footnotehyper}}{\usepackage{footnote}}
\makesavenoteenv{longtable}
\usepackage{graphicx}
\makeatletter
\newsavebox\pandoc@box
\newcommand*\pandocbounded[1]{% scales image to fit in text height/width
  \sbox\pandoc@box{#1}%
  \Gscale@div\@tempa{\textheight}{\dimexpr\ht\pandoc@box+\dp\pandoc@box\relax}%
  \Gscale@div\@tempb{\linewidth}{\wd\pandoc@box}%
  \ifdim\@tempb\p@<\@tempa\p@\let\@tempa\@tempb\fi% select the smaller of both
  \ifdim\@tempa\p@<\p@\scalebox{\@tempa}{\usebox\pandoc@box}%
  \else\usebox{\pandoc@box}%
  \fi%
}
% Set default figure placement to htbp
\def\fps@figure{htbp}
\makeatother



\ifLuaTeX
\usepackage[bidi=basic]{babel}
\else
\usepackage[bidi=default]{babel}
\fi
% get rid of language-specific shorthands (see #6817):
\let\LanguageShortHands\languageshorthands
\def\languageshorthands#1{}
\ifLuaTeX
  \usepackage[english]{selnolig} % disable illegal ligatures
\fi


\setlength{\emergencystretch}{3em} % prevent overfull lines

\providecommand{\tightlist}{%
  \setlength{\itemsep}{0pt}\setlength{\parskip}{0pt}}



 
\usepackage[]{natbib}
\bibliographystyle{plainnat}


\usepackage[a4paper, total={6in, 8in}]{geometry}
\usepackage{tgpagella}
\usepackage{setspace}
\usepackage{tikz}
\usepackage{booktabs, float, multirow}
\usepackage{pdflscape}
\usepackage{amsmath, amsfonts, amssymb}
\usepackage{enumitem}
\usepackage{titlesec}
\usepackage{titling}
\usepackage{tocloft}

\setstretch{1.5}
\setlist{noitemsep, topsep=4pt, leftmargin=1.5em}
\setlength{\parindent}{2em}
\setlength{\parskip}{0pt}

\setcitestyle{round}

\newcommand{\E}{\mathbb{E}}
\newcommand{\Var}{\mathrm{Var}}
\makeatletter
\@ifpackageloaded{caption}{}{\usepackage{caption}}
\AtBeginDocument{%
\ifdefined\contentsname
  \renewcommand*\contentsname{Table of contents}
\else
  \newcommand\contentsname{Table of contents}
\fi
\ifdefined\listfigurename
  \renewcommand*\listfigurename{List of Figures}
\else
  \newcommand\listfigurename{List of Figures}
\fi
\ifdefined\listtablename
  \renewcommand*\listtablename{List of Tables}
\else
  \newcommand\listtablename{List of Tables}
\fi
\ifdefined\figurename
  \renewcommand*\figurename{Figure}
\else
  \newcommand\figurename{Figure}
\fi
\ifdefined\tablename
  \renewcommand*\tablename{Table}
\else
  \newcommand\tablename{Table}
\fi
}
\@ifpackageloaded{float}{}{\usepackage{float}}
\floatstyle{ruled}
\@ifundefined{c@chapter}{\newfloat{codelisting}{h}{lop}}{\newfloat{codelisting}{h}{lop}[chapter]}
\floatname{codelisting}{Listing}
\newcommand*\listoflistings{\listof{codelisting}{List of Listings}}
\usepackage{amsthm}
\theoremstyle{plain}
\newtheorem{proposition}{Proposition}[section]
\theoremstyle{plain}
\newtheorem{corollary}{Corollary}[section]
\theoremstyle{definition}
\newtheorem{definition}{Definition}[section]
\theoremstyle{remark}
\AtBeginDocument{\renewcommand*{\proofname}{Proof}}
\newtheorem*{remark}{Remark}
\newtheorem*{solution}{Solution}
\newtheorem{refremark}{Remark}[section]
\newtheorem{refsolution}{Solution}[section]
\makeatother
\makeatletter
\makeatother
\makeatletter
\@ifpackageloaded{caption}{}{\usepackage{caption}}
\@ifpackageloaded{subcaption}{}{\usepackage{subcaption}}
\makeatother
\usepackage{bookmark}
\IfFileExists{xurl.sty}{\usepackage{xurl}}{} % add URL line breaks if available
\urlstyle{same}
\hypersetup{
  pdftitle={Model Summary},
  pdfauthor={Anthony Sung},
  pdflang={en},
  colorlinks=true,
  linkcolor={blue},
  filecolor={Maroon},
  citecolor={Blue},
  urlcolor={Blue},
  pdfcreator={LaTeX via pandoc}}


\title{Model Summary}
\author{Anthony Sung}
\date{2026-04-02}
\begin{document}
\maketitle

\renewcommand*\contentsname{Table of contents}
{
\hypersetup{linkcolor=}
\setcounter{tocdepth}{3}
\tableofcontents
}

\section{Introduction}\label{sec-intro}

This paper studies endogenous timing in a two-player Tullock contest
under asymmetric information. Labor and management compete over policy
implementation, but only labor privately observes its own valuation of
winning. Before effort choices are made, move order matters: moving
first commits, while moving second preserves flexibility and allows a
response to the opponent. The core questions are: when one side has
private information, who wants to move first, who prefers to wait, and
how is information revealed in equilibrium?

The problem is broadly relevant because many strategic environments
combine timing and asymmetric information, including labor conflict,
electoral competition, lobbying, litigation, and rent seeking. One side
often knows more about how much it values victory. This not only changes
direct effort incentives, but also changes the value of commitment
versus the value of waiting. A first mover may reveal information
through action; a second mover can condition on observed behavior.
Endogenous timing in contests is therefore both a commitment problem and
an information-transmission problem.

The natural benchmark is \citet{fu2006}, which studies endogenous timing
with asymmetric information in a two-player contest. We keep a
parsimonious structure but generalize the model so labor's private value
can follow an arbitrary distribution, and we use binary and uniform
distributions as baseline functional examples. The binary case centers
on mean \(\mu\) with high/low realizations and dispersion controlled by
\(\delta\); the uniform case assigns equal density on
\([\mu-\delta,\mu+\delta]\). The first cleanly separates mean and
dispersion effects, while the second shows a smooth cutoff structure
with continuous types.

The paper has three main results. First, the labor-first subgame (LM)
admits a unique separating equilibrium. When labor moves first and
management responds after observing labor effort, each type's
equilibrium effort fully reveals its private value. The reason is
structural: management's continuation best response depends only on
observed effort and known prize value \(V\), not directly on beliefs.
Hence each type has a strict best response in effort, and pooling cannot
be sustained. This result is distribution-free and applies to binary,
uniform, and other distributions.

Second, timing incentives admit an exact type-by-type characterization.
For each type, the net gain from waiting can be written as an interval
condition, and low types have weakly stronger waiting incentives than
high types. Under binary types this yields a three-way classification:
both types prefer moving first, only the low type prefers waiting, or
both types prefer waiting. Under a uniform distribution there is a
cutoff type \(v^*\): types below the cutoff prefer waiting, and types
above prefer moving first.

Third, mean and dispersion must be considered jointly. Comparing only
the average private value \(\mu\) with management's known value \(V\) is
insufficient for timing predictions. What matters is the interaction of
\(V\), \(\mu\), and dispersion \(\delta\), because \(\delta\) jointly
affects beliefs and type distance.

The contribution is analytical sharpness rather than a larger player set
or richer timing protocol. Relative to \citet{fu2006}, the generalized
setup yields exact type-level timing conditions, a strong no-pooling
result in the informed-first subgame, and a clear decomposition of how
mean and dispersion jointly determine timing incentives.

The rest of the paper is organized as follows. Section
Section~\ref{sec-literature} reviews related work and clarifies the
relation to \citet{fu2006} and subsequent studies. Section
Section~\ref{sec-setup} presents the environment and notation. Section
Section~\ref{sec-equilibria} unifies the SS, ML, and LM equilibria.
Section Section~\ref{sec-timing} characterizes type-level timing
incentives and endogenous timing. Section Section~\ref{sec-welfare}
compares welfare and summarizes the main results in tables.

\section{Literature Review}\label{sec-literature}

\subsection{Main Implications of This
Model}\label{main-implications-of-this-model}

\citet{fu2006} is the closest benchmark. Relative to that paper, our
setup delivers a sharper characterization of timing incentives under
asymmetric information. In the LM subgame, labor (the informed side)
chooses effort first and management responds after observing effort.
Because management's best response depends only on observed effort and
known prize value \(V\), off-path beliefs do not affect continuation
behavior. This eliminates room for belief-supported pooling and yields a
unique separating equilibrium in LM.

Using binary and uniform distributions as functional examples, we
separate the marginal roles of mean private value \(\mu\) and
information range \(\delta\). Timing incentives are not pinned down by
comparing only \(\mu\) and \(V\); the key object is the interval
condition in Proposition~\ref{prp-exact-condition}. Under binary types
this implies exactly three feasible timing profiles, while under a
uniform distribution it implies a cutoff-type structure.

\subsection{Relation to Endogenous-Timing
Literature}\label{relation-to-endogenous-timing-literature}

Classic endogenous-timing results in asymmetric contests suggest weak
players tend to move first and strong players tend to move second. This
logic traces back to the broader framework in \citet{hamilton1990} and
appears in the contest analysis of \citet{fu2006}. More recent work
refines this benchmark. \citet{hoffmann2012} shows that with generalized
success functions and endogenous prizes, simultaneous play can re-emerge
and leadership need not be determined solely by relative strength.
\citet{baikleelee2022} shows that in three-player Tullock contests,
moderate asymmetry can make simultaneous early moves the unique
equilibrium. In a broader sequential all-pay setting,
\citet{dengfuwuzhu2021} establishes robust second-mover payoff
advantages in multi-period structures.

Against this background, our paper provides a tractable two-player
benchmark with endogenous timing and one-sided private information. The
restricted environment permits closed-form type-by-type timing
conditions rather than only qualitative leadership rankings.

\subsection{Relation to Asymmetric-Information and Signaling
Literature}\label{relation-to-asymmetric-information-and-signaling-literature}

This paper also connects to incomplete-information contest work.
\citet{hurley1998}, \citet{warneryd2003}, and \citet{warneryd2012} study
simultaneous contests with uncertainty and show how private/common value
uncertainty changes effort incentives. In this line, \citet{fu2006}
remains the closest reference because informed effort acts as a signal.
Our difference is stronger identification: in our setup, effort
signaling does not merely allow separation; it makes separation unique
in LM.

Other papers introduce explicit communication or signaling stages before
the contest. \citet{fu2013} shows that public messages can play both
commitment and signaling roles; \citet{dentermorgansisak2022} shows that
costly pre-contest signals reveal capability only for sufficiently
strong types. By contrast, our model uses no extra communication
technology: equilibrium effort itself is the only signal. This keeps
comparability with \citet{fu2006} while delivering a stronger no-pooling
result.

\subsection{Relation to Recent
Extensions}\label{relation-to-recent-extensions}

Recent studies also show that richer frictions can reverse standard
timing predictions. \citet{protopappas2023} studies how players
manipulate institutional rules that determine move order in sequential
contests. \citet{protopappas2025} shows that under future cost
uncertainty and asymmetric information, even strong players may
optimally move first. In conflict settings, \citet{dentersisak2015}
shows that incomplete information weakens deterrence because defenders
optimally underinvest when opponent resolve is uncertain.

A key takeaway is that timing predictions are highly sensitive to which
state variable is private information: valuation, future cost, resolve,
or institutional options. Our contribution is to isolate the
private-valuation channel and show that even a minimal information
structure can generate rich timing regions with clear type-level
ordering of waiting incentives.

\subsection{Bottom Line}\label{bottom-line}

The literature comparison suggests this paper is best interpreted as an
analytically sharper extension of \citet{fu2006}. Relative to the
broader contest literature, the central value is not more players or a
richer timing protocol, but sharper identification of when private
information generates first-mover versus second-mover incentives.

\section{Model Setup}\label{sec-setup}

\subsection{Environment}\label{environment}

There are two players: labor (\(L\)) and management (\(M\)). They
compete over which side's preferred policy is implemented. The outcome
is winner-take-all: the winner receives its own prize value, while the
loser receives 0. Contest technology follows \citet{tullock1980};
endogenous timing follows \citet{hamilton1990} and \citet{fu2006}.

\begin{definition}[Asymmetric Prize
Values]\protect\hypertarget{def-prizes}{}\label{def-prizes}

If labor wins, labor obtains \(V_L>0\) and management obtains 0. If
management wins, management obtains \(V_M\equiv V>0\) and labor obtains
0. Prize values need not be symmetric.

\end{definition}

\subsection{Information Structure}\label{information-structure}

Before timing or effort choices, labor privately observes its own prize
value \(V_L\). Management knows only the distribution of \(V_L\). We
assume \(V_L\) has support set \([\mu-\delta,\mu+\delta]\), with
\(\mu>\delta>0\), so that \(\E[V_L]=\mu\).

Define

\begin{equation}\phantomsection\label{eq-u-types}{
u_\ell \equiv \sqrt{\mu-\delta}, \qquad
u_H \equiv \sqrt{\mu+\delta}.
}\end{equation}

We use two baseline distributions.

\begin{definition}[Binary
Distribution]\protect\hypertarget{def-binary}{}\label{def-binary}

\[
V_L \in \{V_\ell,V_H\}=\{\mu-\delta,\mu+\delta\},
\]

with probability \(1/2\) on each type. Hence \(\Var(V_L)=\delta^2\).

\end{definition}

\begin{definition}[Uniform
Distribution]\protect\hypertarget{def-uniform}{}\label{def-uniform}

\[
V_L \sim U[\mu-\delta,\mu+\delta],
\]

with density \(f(v)=1/(2\delta)\). Hence \(\Var(V_L)=\delta^2/3\).

\end{definition}

Both distributions share the same mean \(\mu\) and support set
\([\mu-\delta,\mu+\delta]\), but differ in variance. Equilibrium
expressions depend on two key moments:

\begin{equation}\phantomsection\label{eq-moments}{
\kappa \equiv \E\!\left[\frac{1}{\sqrt{V_L}}\right],
\qquad
\nu \equiv \E\!\left[\sqrt{V_L}\right].
}\end{equation}

\begin{proposition}[Moment
Comparison]\protect\hypertarget{prp-moments}{}\label{prp-moments}

The key moments are:

\begin{longtable}[]{@{}
  >{\raggedright\arraybackslash}p{(\linewidth - 4\tabcolsep) * \real{0.3333}}
  >{\raggedright\arraybackslash}p{(\linewidth - 4\tabcolsep) * \real{0.3333}}
  >{\raggedright\arraybackslash}p{(\linewidth - 4\tabcolsep) * \real{0.3333}}@{}}
\caption{Key moments under binary and uniform
distributions}\label{tbl-moments}\tabularnewline
\toprule\noalign{}
\begin{minipage}[b]{\linewidth}\raggedright
\end{minipage} & \begin{minipage}[b]{\linewidth}\raggedright
Binary
\end{minipage} & \begin{minipage}[b]{\linewidth}\raggedright
Uniform
\end{minipage} \\
\midrule\noalign{}
\endfirsthead
\toprule\noalign{}
\begin{minipage}[b]{\linewidth}\raggedright
\end{minipage} & \begin{minipage}[b]{\linewidth}\raggedright
Binary
\end{minipage} & \begin{minipage}[b]{\linewidth}\raggedright
Uniform
\end{minipage} \\
\midrule\noalign{}
\endhead
\bottomrule\noalign{}
\endlastfoot
\(\Var(V_L)\) & \(\delta^2\) & \(\delta^2/3\) \\
\(\kappa\) & \(\dfrac{1}{2u_\ell}+\dfrac{1}{2u_H}\) &
\(\dfrac{u_H-u_\ell}{\delta}\) \\
\(\nu\) & \(\dfrac{u_\ell+u_H}{2}\) &
\(\dfrac{u_H^3-u_\ell^3}{3\delta}\) \\
\end{longtable}

\end{proposition}

\begin{proof}
For the binary case, moments follow directly from definition. For the
uniform case:

\[
\kappa^U
=\frac{1}{2\delta}\int_{\mu-\delta}^{\mu+\delta} v^{-1/2}\,dv
=\frac{[2\sqrt v]_{\mu-\delta}^{\mu+\delta}}{2\delta}
=\frac{u_H-u_\ell}{\delta},
\]

\[
\nu^U
=\frac{1}{2\delta}\int_{\mu-\delta}^{\mu+\delta} v^{1/2}\,dv
=\frac{[\tfrac{2}{3}v^{3/2}]_{\mu-\delta}^{\mu+\delta}}{2\delta}
=\frac{u_H^3-u_\ell^3}{3\delta}.
\]
\end{proof}

\subsection{Contest Payoffs and Best
Responses}\label{contest-payoffs-and-best-responses}

The Tullock success function is \(p_i=x_i/(x_L+x_M)\). Net expected
payoffs are

\begin{equation}\phantomsection\label{eq-ul}{
U_L(x_L,x_M;V_L)=\frac{x_L}{x_L+x_M}V_L-x_L,
}\end{equation}

\begin{equation}\phantomsection\label{eq-um}{
U_M(x_L,x_M;V)=\frac{x_M}{x_L+x_M}V-x_M.
}\end{equation}

First-order conditions imply best responses:

\begin{equation}\phantomsection\label{eq-brl}{
x_L^*(x_M;V_L)=\sqrt{V_Lx_M}-x_M,
}\end{equation}

\begin{equation}\phantomsection\label{eq-brm}{
x_M^*(x_L;V)=\sqrt{Vx_L}-x_L.
}\end{equation}

Management's best response Equation~\ref{eq-brm} depends only on
observed effort \(x_L\) and known value \(V\). Beliefs about labor types
do not directly enter the second-stage best response. This is central
for the separation-versus-pooling result later.

\section{Equilibrium Analysis}\label{sec-equilibria}

This section unifies the three subgames: SS (simultaneous moves), ML
(management moves first), and LM (labor moves first). For comparison, we
use the best responses in Equation~\ref{eq-brl} and
Equation~\ref{eq-brm} throughout.

\subsection{SS: Simultaneous Moves}\label{sec-ss}

Both players choose effort simultaneously. Labor chooses \(x_L\) by type
\(V_L\), while management chooses \(x_M\) under incomplete information
about \(V_L\). Substituting Equation~\ref{eq-brl} into management's
objective and taking expectation over types gives

\[
\E[U_M] = V\sqrt{x_M}\,\kappa - x_M,
\qquad
\frac{\partial \E[U_M]}{\partial x_M}
  = \frac{V\kappa}{2\sqrt{x_M}} - 1.
\]

\begin{proposition}[Simultaneous-Move
Equilibrium]\protect\hypertarget{prp-ss}{}\label{prp-ss}

The SS equilibrium is

\begin{equation}\phantomsection\label{eq-ss-efforts}{
x_M^{\mathrm{SS}} = \frac{V^2\kappa^2}{4}, \qquad
x_L^{\mathrm{SS}}(V_L)
  = \frac{V\kappa}{2}\sqrt{V_L} - \frac{V^2\kappa^2}{4}.
}\end{equation}

Equilibrium payoffs are

\begin{equation}\phantomsection\label{eq-ul-ss}{
U_L^{\mathrm{SS}}(V_L)
  = \left(\sqrt{V_L} - \frac{V\kappa}{2}\right)^2,
}\end{equation}

\begin{equation}\phantomsection\label{eq-um-ss}{
U_M^{\mathrm{SS}}(V_L)
  = \frac{V^2\kappa}{2\sqrt{V_L}} - \frac{V^2\kappa^2}{4}.
}\end{equation}

Expected payoffs are

\begin{equation}\phantomsection\label{eq-eul-ss}{
\E\!\left[U_L^{\mathrm{SS}}\right]
  = \mu - V\kappa\nu + \frac{V^2\kappa^2}{4},
}\end{equation}

\begin{equation}\phantomsection\label{eq-eum-ss}{
\E\!\left[U_M^{\mathrm{SS}}\right]
  = \frac{V^2\kappa^2}{4}.
}\end{equation}

\end{proposition}

\begin{proof}
From Equation~\ref{eq-brl}, labor's best response for type \(v\) is
\(x_L^*(x_M;v)=\sqrt{vx_M}-x_M\). Substitute into Equation~\ref{eq-um}:
\[
U_M(x_L^*,x_M;V)
=\frac{x_M}{x_M+\sqrt{vx_M}-x_M}V-x_M
=V\sqrt{\frac{x_M}{v}}-x_M.
\] Taking expectation over \(v\): \[
\E[U_M]=V\sqrt{x_M}\,\E[v^{-1/2}]-x_M
=V\sqrt{x_M}\,\kappa-x_M.
\] FOC and SOC are \[
\frac{\partial \E[U_M]}{\partial x_M}
=\frac{V\kappa}{2\sqrt{x_M}}-1=0,
\qquad
\frac{\partial^2 \E[U_M]}{\partial x_M^2}
=-\frac{V\kappa}{4x_M^{3/2}}<0,
\] so the unique maximizer is \(x_M^{\mathrm{SS}}=V^2\kappa^2/4\). Then
Equation~\ref{eq-brl} gives \[
x_L^{\mathrm{SS}}(v)
=\sqrt{v\cdot\frac{V^2\kappa^2}{4}}-\frac{V^2\kappa^2}{4}
=\frac{V\kappa}{2}\sqrt v-\frac{V^2\kappa^2}{4}.
\] Substitute into Equation~\ref{eq-ul} and Equation~\ref{eq-um}: \[
U_L^{\mathrm{SS}}(v)
=\frac{x_L^{\mathrm{SS}}}{x_L^{\mathrm{SS}}+x_M^{\mathrm{SS}}}v-x_L^{\mathrm{SS}}
=\left(\sqrt v-\frac{V\kappa}{2}\right)^2,
\] \[
U_M^{\mathrm{SS}}(v)
=\frac{x_M^{\mathrm{SS}}}{x_L^{\mathrm{SS}}+x_M^{\mathrm{SS}}}V-x_M^{\mathrm{SS}}
=\frac{V^2\kappa}{2\sqrt v}-\frac{V^2\kappa^2}{4}.
\] Finally, \[
\E[U_L^{\mathrm{SS}}]
=\E\!\left[v-V\kappa\sqrt v+\frac{V^2\kappa^2}{4}\right]
=\mu-V\kappa\nu+\frac{V^2\kappa^2}{4},
\] \[
\E[U_M^{\mathrm{SS}}]
=\E\!\left[\frac{V^2\kappa}{2\sqrt v}-\frac{V^2\kappa^2}{4}\right]
=\frac{V^2\kappa}{2}\E[v^{-1/2}]-\frac{V^2\kappa^2}{4}
=\frac{V^2\kappa^2}{4}.
\]
\end{proof}

Substituting moments from Table~\ref{tbl-moments} into
Proposition~\ref{prp-ss} gives distribution-specific formulas.

\textbf{Binary.} Let \(\kappa^B = \frac{1}{2u_\ell} + \frac{1}{2u_H}\)
and \(\nu^B = \frac{u_\ell + u_H}{2}\):

\[
x_M^{\mathrm{SS},B} = \frac{V^2(\kappa^B)^2}{4},
\qquad
\E\!\left[U_L^{\mathrm{SS},B}\right]
  = \mu - V\kappa^B\nu^B + \frac{V^2(\kappa^B)^2}{4}.
\]

\textbf{Uniform.} Let \(\kappa^U = \frac{u_H - u_\ell}{\delta}\) and
\(\nu^U = \frac{u_H^3 - u_\ell^3}{3\delta}\):

\[
x_M^{\mathrm{SS},U} = \frac{V^2(\kappa^U)^2}{4},
\qquad
\E\!\left[U_L^{\mathrm{SS},U}\right]
  = \mu - V\kappa^U\nu^U + \frac{V^2(\kappa^U)^2}{4}.
\]

Because \(v\mapsto v^{-1/2}\) is convex, binary types imply larger
\(\kappa\) than uniform (Jensen), which increases management effort in
SS.

\subsection{ML: Management Moves First}\label{sec-ml}

Management commits to \(x_M\) before observing labor's type. Labor then
responds type-by-type using Equation~\ref{eq-brl}. Hence management
solves the same expected objective as in SS.

\begin{proposition}[ML
Equilibrium]\protect\hypertarget{prp-ml}{}\label{prp-ml}

\begin{equation}\phantomsection\label{eq-ml-efforts}{
x_M^{\mathrm{ML}} = \frac{V^2\kappa^2}{4}, \qquad
x_L^{\mathrm{ML}}(V_L)
  = \frac{V\kappa}{2}\sqrt{V_L} - \frac{V^2\kappa^2}{4}.
}\end{equation}

Payoffs coincide with Equation~\ref{eq-ul-ss}--Equation~\ref{eq-eum-ss}:

\begin{equation}\phantomsection\label{eq-eul-ml}{
\E\!\left[U_L^{\mathrm{ML}}\right]
  = \mu - V\kappa\nu + \frac{V^2\kappa^2}{4},
}\end{equation}

\begin{equation}\phantomsection\label{eq-eum-ml}{
\E\!\left[U_M^{\mathrm{ML}}\right]
  = \frac{V^2\kappa^2}{4}.
}\end{equation}

\end{proposition}

\begin{proof}
In ML, stage 2 still gives labor's type-contingent response
\(x_L^*(x_M;v)=\sqrt{vx_M}-x_M\). Therefore stage-1 management solves \[
\max_{x_M\ge 0}\;\E\!\left[V\sqrt{\frac{x_M}{V_L}}-x_M\right]
=\max_{x_M\ge 0}\;\left(V\sqrt{x_M}\,\kappa-x_M\right),
\] which is exactly the SS problem. The same FOC/SOC imply the same
optimal \(x_M\), and Equation~\ref{eq-brl} then implies the same
\(x_L(v)\). Hence all payoffs coincide with SS.
\end{proof}

Distribution-specific expressions are therefore identical to the SS
formulas above.

\subsection{LM: Labor Moves First}\label{sec-lm}

Labor moves in period 1 and management responds in period 2. Observing
\(x_L\), management uses Equation~\ref{eq-brm}. For labor type
\(V_L=v\):

\[
\max_{x_L \geq 0}\;
W(x_L;v)
\equiv
\frac{x_L}{x_L + x_M^*(x_L)}v - x_L
=
\frac{v}{\sqrt{V}}\sqrt{x_L} - x_L.
\]

\(W(x_L;v)\) is strictly concave in \(x_L\), so each type has a unique
optimizer.

\begin{proposition}[Separating Effort
Levels]\protect\hypertarget{prp-separating}{}\label{prp-separating}

For type \(v\),

\begin{equation}\phantomsection\label{eq-lm-efforts}{
x_L^{\mathrm{LM}}(v) = \frac{v^2}{4V},
\qquad
x_M^{\mathrm{LM}}(v) = \frac{v(2V - v)}{4V},
}\end{equation}

\begin{equation}\phantomsection\label{eq-ul-lm}{
U_L^{\mathrm{LM}}(v) = \frac{v^2}{4V},
}\end{equation}

\begin{equation}\phantomsection\label{eq-um-lm}{
U_M^{\mathrm{LM}}(v) = \frac{(2V - v)^2}{4V}.
}\end{equation}

\end{proposition}

\begin{proof}
Labor solves \[
\max_{x_L\ge 0}\;W(x_L;v)=\max_{x_L\ge 0}\left(\frac{v}{\sqrt V}\sqrt{x_L}-x_L\right).
\] FOC: \[
\frac{v}{2\sqrt{Vx_L}} - 1 = 0
\quad\Rightarrow\quad
x_L^{\mathrm{LM}}(v)=\frac{v^2}{4V}.
\] SOC: \[
\frac{\partial^2 W}{\partial x_L^2}=-\frac{v}{4\sqrt V}x_L^{-3/2}<0,
\] so this is the unique global optimum. Then Equation~\ref{eq-brm}
implies \[
x_M^{\mathrm{LM}}(v)
=\sqrt{V\cdot\frac{v^2}{4V}}-\frac{v^2}{4V}
=\frac{v(2V-v)}{4V}.
\] Substitute into Equation~\ref{eq-ul} and Equation~\ref{eq-um} to get
Equation~\ref{eq-ul-lm} and Equation~\ref{eq-um-lm}.
\end{proof}

\begin{proposition}[No Pooling
Equilibrium]\protect\hypertarget{prp-no-pool}{}\label{prp-no-pool}

The LM subgame admits a unique PBE, and it is separating.

\end{proposition}

\begin{proof}
Management's continuation best response Equation~\ref{eq-brm} depends
only on observed \(x_L\) and \(V\), not on beliefs about labor's type.
Hence each type \(v\) solves the same standalone problem: \[
W(x_L;v)=\frac{v}{\sqrt V}\sqrt{x_L}-x_L,
\] with unique maximizer \(x_L^{\mathrm{LM}}(v)=v^2/(4V)\). Therefore
any profile prescribing a different effort for type \(v\) is strictly
deviable.

Moreover, \[
\frac{d}{dv}x_L^{\mathrm{LM}}(v)=\frac{v}{2V}>0,
\] so different types choose different efforts. Pooling cannot satisfy
type-wise optimality, so only separating behavior survives in
equilibrium.
\end{proof}

This argument is distribution-free: binary, uniform, or any continuous
type distribution all lead to the same uniqueness logic in LM.

\subsubsection{Expected Payoffs}\label{expected-payoffs}

Under the unique separating equilibrium,

\begin{equation}\phantomsection\label{eq-eul-lm}{
\E\!\left[U_L^{\mathrm{LM}}\right]
  = \frac{\E[V_L^2]}{4V}
  = \frac{\mu^2 + \sigma^2}{4V},
}\end{equation}

\begin{equation}\phantomsection\label{eq-eum-lm}{
\E\!\left[U_M^{\mathrm{LM}}\right]
  = \frac{4V^2 - 4V\mu + \mu^2 + \sigma^2}{4V},
}\end{equation}

where \(\sigma^2=\Var(V_L)\).

\begin{proof}
From Equation~\ref{eq-ul-lm}, \(U_L^{\mathrm{LM}}(v)=v^2/(4V)\), so \[
\E[U_L^{\mathrm{LM}}]
=\E\!\left[\frac{V_L^2}{4V}\right]
=\frac{\E[V_L^2]}{4V}
=\frac{\mu^2+\sigma^2}{4V}.
\] From Equation~\ref{eq-um-lm}, \[
\E[U_M^{\mathrm{LM}}]
=\E\!\left[\frac{(2V-V_L)^2}{4V}\right]
=\frac{4V^2-4V\E[V_L]+\E[V_L^2]}{4V}
=\frac{4V^2-4V\mu+\mu^2+\sigma^2}{4V}.
\]
\end{proof}

\subsubsection{Functional Examples}\label{sec-lm-examples}

\textbf{Binary} (\(\sigma^2=\delta^2\)):

\[
\E\!\left[U_L^{\mathrm{LM},B}\right]
  = \frac{\mu^2 + \delta^2}{4V},
\qquad
\E\!\left[U_M^{\mathrm{LM},B}\right]
  = \frac{4V^2 - 4V\mu + \mu^2 + \delta^2}{4V}.
\]

\textbf{Uniform} (\(\sigma^2=\delta^2/3\)):

\[
\E\!\left[U_L^{\mathrm{LM},U}\right]
  = \frac{3\mu^2 + \delta^2}{12V},
\qquad
\E\!\left[U_M^{\mathrm{LM},U}\right]
  = \frac{12V^2 - 12V\mu + 3\mu^2 + \delta^2}{12V}.
\]

Since \(\delta^2>\delta^2/3\), holding \((\mu,\delta,V)\) fixed, labor's
LM expected payoff is higher under binary than uniform types.

\subsection{Timing Incentives and Endogenous Timing}\label{sec-timing}

\subsection{Type-Level Incentives}\label{type-level-incentives}

For type \(v\), define the gain from waiting (second move) relative to
moving first:

\begin{equation}\phantomsection\label{eq-rent-function}{
R(v) \equiv U_L^{\mathrm{SS}}(v) - U_L^{\mathrm{LM}}(v)
  = \left(\sqrt{v} - \frac{V\kappa}{2}\right)^2 - \frac{v^2}{4V}.
}\end{equation}

Thus, \(R(v)>0\) means type \(v\) prefers waiting (SS/ML), while
\(R(v)<0\) means it prefers moving first (LM). Let \(u=\sqrt v\) and
\(c\equiv V\kappa/2\); then

\begin{equation}\phantomsection\label{eq-rent-factored}{
R(u)
=
\left(c-u-\frac{u^2}{2\sqrt V}\right)
\left(c-u+\frac{u^2}{2\sqrt V}\right).
}\end{equation}

\begin{proposition}[Exact Type
Condition]\protect\hypertarget{prp-exact-condition}{}\label{prp-exact-condition}

Let \(c\equiv V\kappa/2\). Type \(v\) prefers moving first if and only
if \(c\) lies in

\[
I(v)
\equiv
\left[\sqrt v-\frac{v}{2\sqrt V},\;\sqrt v+\frac{v}{2\sqrt V}\right].
\]

Equivalently, type \(v\) prefers waiting when \(c\) lies strictly
outside \(I(v)\).

\end{proposition}

\begin{proof}
From \[
R(v)=\left(\sqrt v-c\right)^2-\frac{v^2}{4V}\le 0
\iff
\left|\sqrt v-c\right|\le \frac{v}{2\sqrt V},
\] we obtain \[
-\frac{v}{2\sqrt V}\le \sqrt v-c\le \frac{v}{2\sqrt V}
\iff
\sqrt v-\frac{v}{2\sqrt V}\le c\le \sqrt v+\frac{v}{2\sqrt V},
\] which is exactly \(c\in I(v)\).
\end{proof}

\begin{proposition}[Low Types Have Stronger Waiting
Incentives]\protect\hypertarget{prp-low-more-wait}{}\label{prp-low-more-wait}

For any \(v_1<v_2\),

\[
R(v_1)\ge R(v_2),
\]

with strict inequality when \(v_1\neq v_2\).

\end{proposition}

\begin{proof}
Write \(R(v)=v-2c\sqrt v+c^2-v^2/(4V)\). The linear term in \(v\) and
the concavity of \(\sqrt v\) increase waiting value for low types, while
the quadratic term \(v^2/(4V)\) raises first-mover value faster for high
types. Therefore \(R(v)\) is decreasing in \(v\) on the relevant
support, implying the stated ordering.
\end{proof}

\subsection{Distribution-Specific Timing
Structures}\label{distribution-specific-timing-structures}

\begin{corollary}[Binary Case: Exactly Three
Profiles]\protect\hypertarget{cor-three-profiles}{}\label{cor-three-profiles}

Under binary types (\(V_L\in\{V_\ell,V_H\}\)), only three patterns are
possible:

\begin{enumerate}
\def\labelenumi{\arabic{enumi}.}
\tightlist
\item
  Both types prefer moving first: \(R(V_\ell)\le 0\).
\item
  Only the low type prefers waiting: \(R(V_\ell)>0\ge R(V_H)\).
\item
  Both types prefer waiting: \(R(V_H)>0\).
\end{enumerate}

The reverse pattern (only the high type prefers waiting) is impossible
by Proposition~\ref{prp-low-more-wait}.

\end{corollary}

\begin{proposition}[Uniform Case: Cutoff
Type]\protect\hypertarget{prp-uniform-cutoff}{}\label{prp-uniform-cutoff}

Under a uniform distribution, \(R(v)\) is continuous and decreasing by
Proposition~\ref{prp-low-more-wait}. Hence there exists a cutoff type

\[
v^*\in[\mu-\delta,\mu+\delta]\cup\{+\infty,-\infty\},
\]

such that:

\begin{itemize}
\tightlist
\item
  if \(v^*\in(\mu-\delta,\mu+\delta)\): types \(v<v^*\) prefer waiting,
  types \(v>v^*\) prefer moving first;
\item
  if \(v^*\le \mu-\delta\): all types prefer moving first;
\item
  if \(v^*\ge \mu+\delta\): all types prefer waiting.
\end{itemize}

The cutoff is defined implicitly by \(R(v^*)=0\).

\end{proposition}

\begin{proof}
\(R(v)\) is continuous on the support interval and monotone decreasing.
If \(R(\mu-\delta)\le 0\), then all types prefer moving first. If
\(R(\mu+\delta)>0\), then all types prefer waiting. Otherwise,
\(R(\mu-\delta)>0\ge R(\mu+\delta)\), so the intermediate value theorem
gives at least one root and monotonicity gives uniqueness.
\end{proof}

Comparing only \(\mu\) and \(V\) is not enough. The relevant object is
\(c=V\kappa/2\), and \(\kappa\) depends jointly on mean and dispersion.
Therefore \(\mu<V\) does not rule out waiting incentives, and \(\mu>V\)
does not guarantee them.

\subsection{Management Incentives and Endogenous
Timing}\label{sec-endogenous}

\begin{proposition}[Management's Timing
Preference]\protect\hypertarget{prp-mgmt}{}\label{prp-mgmt}

Management prefers moving first (ML) to moving second (LM) if and only
if

\begin{equation}\phantomsection\label{eq-mgmt-cond}{
\frac{V^2\kappa^2}{4} \ge \frac{4V^2-4V\mu+\mu^2+\sigma^2}{4V}.
}\end{equation}

\end{proposition}

\begin{proof}
From Equation~\ref{eq-eum-ml} and Equation~\ref{eq-eum-lm}: \[
\E\!\left[U_M^{\mathrm{ML}}\right]=\frac{V^2\kappa^2}{4},
\qquad
\E\!\left[U_M^{\mathrm{LM}}\right]=\frac{4V^2-4V\mu+\mu^2+\sigma^2}{4V}.
\] The condition \(\E[U_M^{\mathrm{ML}}]\ge\E[U_M^{\mathrm{LM}}]\) is
exactly Equation~\ref{eq-mgmt-cond}.
\end{proof}

\subsubsection{Functional Examples}\label{functional-examples}

\textbf{Binary} (\(\sigma^2=\delta^2\)):

\[
V^3(\kappa^B)^2 \ge 4V^2-4V\mu+\mu^2+\delta^2.
\]

\textbf{Uniform} (\(\sigma^2=\delta^2/3\)):

\[
V^3(\kappa^U)^2 \ge 4V^2-4V\mu+\mu^2+\frac{\delta^2}{3}.
\]

\subsubsection{Pooling Ambiguity}\label{pooling-ambiguity}

Because LM is uniquely separating (Proposition~\ref{prp-no-pool}),
endogenous timing can vary by type without pooling-equilibrium
ambiguity.

\subsection{Welfare and Result Tables}\label{sec-welfare}

\subsection{Welfare Comparison}\label{welfare-comparison}

Total expected welfare is \(\E[W]=\E[U_L]+\E[U_M]\).

\begin{equation}\phantomsection\label{eq-ew-ss}{
\E[W^{\mathrm{SS}}]=\E[W^{\mathrm{ML}}]
=\mu-V\kappa\nu+\frac{V^2\kappa^2}{2},
}\end{equation}

\begin{equation}\phantomsection\label{eq-ew-lm}{
\E[W^{\mathrm{LM}}]
=V-\mu+\frac{\mu^2+\sigma^2}{2V}.
}\end{equation}

\begin{proof}
By definition, \(W=U_L+U_M\). For SS (and ML, which is
payoff-equivalent): \[
\E[W^{\mathrm{SS}}]
=\E[U_L^{\mathrm{SS}}]+\E[U_M^{\mathrm{SS}}]
=\left(\mu-V\kappa\nu+\frac{V^2\kappa^2}{4}\right)+\frac{V^2\kappa^2}{4}
=\mu-V\kappa\nu+\frac{V^2\kappa^2}{2}.
\] For LM: \[
\E[W^{\mathrm{LM}}]
=\E[U_L^{\mathrm{LM}}]+\E[U_M^{\mathrm{LM}}]
=\frac{\mu^2+\sigma^2}{4V}+\frac{4V^2-4V\mu+\mu^2+\sigma^2}{4V}
=V-\mu+\frac{\mu^2+\sigma^2}{2V}.
\]
\end{proof}

The sign of \(\E[W^{\mathrm{SS}}]-\E[W^{\mathrm{LM}}]\) depends jointly
on \(V\), \(\mu\), \(\delta\), and distributional shape.

\textbf{Binary} (\(\sigma^2=\delta^2\), \(\kappa=\kappa^B\),
\(\nu=\nu^B\)):

\[
\E[W^{\mathrm{SS},B}] = \mu - V\kappa^B\nu^B + \frac{V^2(\kappa^B)^2}{2},
\qquad
\E[W^{\mathrm{LM},B}] = V - \mu + \frac{\mu^2 + \delta^2}{2V}.
\]

\textbf{Uniform} (\(\sigma^2=\delta^2/3\), \(\kappa=\kappa^U\),
\(\nu=\nu^U\)):

\[
\E[W^{\mathrm{SS},U}] = \mu - V\kappa^U\nu^U + \frac{V^2(\kappa^U)^2}{2},
\qquad
\E[W^{\mathrm{LM},U}] = V - \mu + \frac{3\mu^2 + \delta^2}{6V}.
\]

\subsection{Result Tables}\label{result-tables}

\begin{longtable}[]{@{}
  >{\raggedright\arraybackslash}p{(\linewidth - 6\tabcolsep) * \real{0.2500}}
  >{\raggedright\arraybackslash}p{(\linewidth - 6\tabcolsep) * \real{0.2500}}
  >{\raggedright\arraybackslash}p{(\linewidth - 6\tabcolsep) * \real{0.2500}}
  >{\raggedright\arraybackslash}p{(\linewidth - 6\tabcolsep) * \real{0.2500}}@{}}
\caption{Equilibrium summary in the generalized
model}\label{tbl-summary-general}\tabularnewline
\toprule\noalign{}
\begin{minipage}[b]{\linewidth}\raggedright
Object
\end{minipage} & \begin{minipage}[b]{\linewidth}\raggedright
SS (Simultaneous)
\end{minipage} & \begin{minipage}[b]{\linewidth}\raggedright
ML (Management First)
\end{minipage} & \begin{minipage}[b]{\linewidth}\raggedright
LM (Labor First)
\end{minipage} \\
\midrule\noalign{}
\endfirsthead
\toprule\noalign{}
\begin{minipage}[b]{\linewidth}\raggedright
Object
\end{minipage} & \begin{minipage}[b]{\linewidth}\raggedright
SS (Simultaneous)
\end{minipage} & \begin{minipage}[b]{\linewidth}\raggedright
ML (Management First)
\end{minipage} & \begin{minipage}[b]{\linewidth}\raggedright
LM (Labor First)
\end{minipage} \\
\midrule\noalign{}
\endhead
\bottomrule\noalign{}
\endlastfoot
Management information & Prior only & Prior only & Observes \(x_L\) and
infers type \\
Equilibrium type & Bayesian Nash & Same as SS & Unique separating PBE \\
\(x_M\) & \(\dfrac{V^2\kappa^2}{4}\) & \(\dfrac{V^2\kappa^2}{4}\) &
\(\dfrac{v(2V-v)}{4V}\) \\
\(x_L\) & \(\dfrac{V\kappa}{2}\sqrt{v} - \dfrac{V^2\kappa^2}{4}\) & Same
as SS & \(\dfrac{v^2}{4V}\) \\
\(U_L\) & \(\left(\sqrt{v}-\dfrac{V\kappa}{2}\right)^{\!2}\) & Same as
SS & \(\dfrac{v^2}{4V}\) \\
\(U_M\) & \(\dfrac{V^2\kappa}{2\sqrt{v}} - \dfrac{V^2\kappa^2}{4}\) &
Same as SS & \(\dfrac{(2V-v)^2}{4V}\) \\
\(\E[U_L]\) & \(\mu - V\kappa\nu + \dfrac{V^2\kappa^2}{4}\) & Same as SS
& \(\dfrac{\mu^2+\sigma^2}{4V}\) \\
\(\E[U_M]\) & \(\dfrac{V^2\kappa^2}{4}\) & \(\dfrac{V^2\kappa^2}{4}\) &
\(\dfrac{4V^2-4V\mu+\mu^2+\sigma^2}{4V}\) \\
\end{longtable}

where \(\kappa=\E[V_L^{-1/2}]\), \(\nu=\E[V_L^{1/2}]\), and
\(\sigma^2=\Var(V_L)\).

\bigskip

\begin{longtable}[]{@{}
  >{\raggedright\arraybackslash}p{(\linewidth - 4\tabcolsep) * \real{0.3333}}
  >{\raggedright\arraybackslash}p{(\linewidth - 4\tabcolsep) * \real{0.3333}}
  >{\raggedright\arraybackslash}p{(\linewidth - 4\tabcolsep) * \real{0.3333}}@{}}
\caption{Moments and timing structure by
distribution}\label{tbl-summary-distributions}\tabularnewline
\toprule\noalign{}
\begin{minipage}[b]{\linewidth}\raggedright
Moment / Structure
\end{minipage} & \begin{minipage}[b]{\linewidth}\raggedright
Binary
\end{minipage} & \begin{minipage}[b]{\linewidth}\raggedright
Uniform
\end{minipage} \\
\midrule\noalign{}
\endfirsthead
\toprule\noalign{}
\begin{minipage}[b]{\linewidth}\raggedright
Moment / Structure
\end{minipage} & \begin{minipage}[b]{\linewidth}\raggedright
Binary
\end{minipage} & \begin{minipage}[b]{\linewidth}\raggedright
Uniform
\end{minipage} \\
\midrule\noalign{}
\endhead
\bottomrule\noalign{}
\endlastfoot
\(\sigma^2\) & \(\delta^2\) & \(\delta^2/3\) \\
\(\kappa\) & \(\dfrac{1}{2u_\ell}+\dfrac{1}{2u_H}\) &
\(\dfrac{u_H-u_\ell}{\delta}\) \\
\(\nu\) & \(\dfrac{u_\ell+u_H}{2}\) &
\(\dfrac{u_H^3-u_\ell^3}{3\delta}\) \\
Timing structure & Three discrete profiles & Cutoff type \(v^*\) \\
\end{longtable}


\bibliography{refs.bib}



\end{document}
